\documentclass[11pt,a4paper]{article}
\usepackage{gtu-el-yazisi}
\begin{document}
\notbilgi{Analog Haberleşme}{ELM362 Final (2007), Soru 1 çözümü}{Z2.pdf}

\sayfa{1}
\begin{ornek}[1]
  \textit{ELM362 Analog Haberleşme Sistemleri Final Sınavı, 8.6.2007.} (30p.) $x(t)=2\,\mathrm{Sinc}(40\pi t)$ işareti, transfer fonksiyonu $H(\omega)$ olan bir kanaldan iletilmektedir. Çıkış işareti $y(t)=20\,\mathrm{Sinc}(40\pi t-200\pi)$ olduğuna göre,
  \begin{enumerate}[label=\alph*),nosep,leftmargin=*]
    \item $H(\omega)$'yı bulunuz ve genlik, faz spektrumlarını çiziniz.
    \item $h(t)$'yi bulunuz.
    \item Bu sistem bozulmasız mıdır? Sistemin girişine $x(t)=t\,\mathrm{Sinc}^2(\pi t)$ işareti uygulanırsa çıkış ne olur?
    \item $x(t)=2\,\mathrm{Sinc}(40\pi t)$ işareti için, çıkışın enerji spektral yoğunluğunu bulunuz.
    \item Çıkış işaretinin enerjisini bulunuz.
  \end{enumerate}
\end{ornek}

a)
\formulkutu{A z\,\operatorname{sinc}(tz) \overset{F}{\longleftrightarrow} A\,\Pi\!\left(\frac{f}{z}\right)}
\[
  x(t)=\underbrace{2}_{Az}\operatorname{sinc}\bigl(\underbrace{40\pi}_{z}\,t\bigr)
  \donusum \frac{1}{20\pi}\,\Pi\!\left(\frac{f}{40\pi}\right)=X(f)
\]
\[
  \begin{aligned}
    Az &= 2\\
    A\cdot 40\pi &= 2\\
    A &= \frac{2}{40\pi}=\frac{1}{20\pi}
  \end{aligned}
\]
\[
  \underbrace{20}_{Az}\operatorname{sinc}\bigl(\underbrace{40\pi}_{z}\,t\bigr)
  \donusum \frac{1}{2\pi}\,\Pi\!\left(\frac{f}{40\pi}\right),
  \qquad Az=20,\quad A=\frac{20}{40\pi}=\frac{1}{2\pi}
\]
\formulkutu[Zamanda öteleme]{x(t-t_0)\overset{F}{\longleftrightarrow}X(f)\,e^{-j2\pi f t_0}}
\[
  y(t)=20\operatorname{sinc}\bigl(40\pi(t-5)\bigr)\donusum
  \frac{1}{2\pi}\,\Pi\!\left(\frac{f}{40\pi}\right)\cdot e^{-j2\pi f 5}=Y(f)
\]
\noindent\begin{minipage}{\linewidth}
\[
  \begin{aligned}
    x(t)\circledast h(t)&=y(t)\\
    X(f)\cdot H(f)&=Y(f)\\
    H(f)=\frac{X(f)}{Y(f)}&=\frac{10\,\dfrac{1}{20\pi}\,\Pi\!\left(\dfrac{f}{40\pi}\right)}{\dfrac{1}{2\pi}\,\Pi\!\left(\dfrac{f}{40\pi}\right)\cdot e^{-j2\pi f 5}}
    =\frac{1}{10}\,\Pi\!\left(\frac{f}{40\pi}\right)e^{j2\pi f 5}\text{\footnotemark}
  \end{aligned}
\]
\footnotetext{Claude notu: muhtemelen $H(f)=Y(f)/X(f)$ olmalı; bu durumda $H(f)=10\,\Pi\!\left(\frac{f}{40\pi}\right)e^{-j2\pi f 5}$ çıkardı.}
\end{minipage}
\noindent\begin{minipage}[t]{0.48\linewidth}
  \centering
  \begin{tikzpicture}
    \begin{axis}[width=6.4cm,height=4.2cm,xmin=-100,xmax=100,ymin=0,ymax=0.17,
        xlabel={$f$},ylabel={$|H(f)|$},xtick={-62.83,62.83},xticklabels={$-20\pi$,$20\pi$},
        ytick={0.1},yticklabels={$\frac{1}{10}$}]
      \addplot[murekkep,line width=1pt] coordinates {(-62.83,0) (-62.83,0.1) (62.83,0.1) (62.83,0)};
      \draw[murekkep,<->] ([yshift=-8mm]axis cs:-62.83,0) -- ([yshift=-8mm]axis cs:62.83,0)
        node[midway,below]{\small $40\pi$};
    \end{axis}
  \end{tikzpicture}
\end{minipage}\hfill
\begin{minipage}[t]{0.48\linewidth}
  \centering
  \begin{tikzpicture}
    \begin{axis}[width=6.4cm,height=4.2cm,xmin=-3,xmax=3,ymin=-40,ymax=40,
        xlabel={$f$},ylabel={$\theta_H(f)$},xtick={1},ytick={31.4159},yticklabels={$10\pi$}]
      \addplot[murekkep,line width=1pt,domain=-2.6:2.6] {31.4159*x};
      \draw[murekkep,dashed] (axis cs:0,31.4159) -- (axis cs:1,31.4159) -- (axis cs:1,0);
    \end{axis}
  \end{tikzpicture}
\end{minipage}

\sayfa{2}
b)
\[
  Az\,\operatorname{sinc}(tz)\overset{F}{\longleftrightarrow}A\,\Pi\!\left(\frac{f}{z}\right),
  \qquad A\rightarrow\frac{1}{10},\quad z\rightarrow 40\pi
\]
\[
  \frac{1}{10}\,40\pi\operatorname{sinc}(40\pi t)\overset{F}{\longleftrightarrow}\frac{1}{10}\,\Pi\!\left(\frac{f}{40\pi}\right)
\]
\[
  \overset{h(t)}{4\pi\operatorname{sinc}\bigl(40\pi(t+5)\bigr)}\overset{F}{\longleftrightarrow}
  \frac{1}{10}\,\Pi\!\left(\frac{f}{40\pi}\right)\cdot e^{j2\pi f 5}=H(f)
\]
c) Sistem her yerde bozulmasız değildir ama $-20\pi$ $20\pi$ arasında bozulmasız gibi davranır çünkü transfer fonksiyonun genliği sabit açısı lineerdir.
\[ x(t)=t\operatorname{sinc}^2(\pi t) \]
\formulkutu[zaman ile çarpma]{t^n x(t)\overset{F}{\longleftrightarrow}\left(\frac{j}{2\pi}\right)^{\!n}\frac{d^n}{df^n}X(f)}
\formulkutu{Az\operatorname{sinc}^2(tz)\overset{F}{\longleftrightarrow}A\,\Lambda\!\left(\frac{f}{z}\right)}
\[
  \operatorname{sinc}^2(\underbrace{\pi}_{z}\,t)\donusum\frac{1}{\pi}\,\Lambda\!\left(\frac{f}{\pi}\right),
  \qquad Az=1,\quad A=\frac{1}{\pi}
\]
\[
  t\operatorname{sinc}^2(\pi t)\donusum\frac{1}{\pi}\left(\frac{j}{2\pi}\right)\frac{d}{df}\,\Lambda\!\left(\frac{f}{\pi}\right)
\]
\begin{center}
  \begin{tikzpicture}
    % ust cizim: Lambda(f/pi)
    \draw[-{Stealth[length=2.2mm]},line width=0.7pt] (-3.5,0) -- (3.8,0) node[below]{$f$};
    \draw[-{Stealth[length=2.2mm]},line width=0.7pt] (0,-0.2) -- (0,2.2);
    \draw[murekkep,line width=1pt] (-2.6,0) -- (0,1.8) -- (2.6,0);
    \node[right] at (0.05,1.8) {$1$};
    \draw[<->] (-2.6,-0.35) -- (0,-0.35) node[midway,below]{$\pi$};
    \draw[<->] (0,-0.35) -- (2.6,-0.35) node[midway,below]{$\pi$};
    \node[anchor=north east] at (-2.6,-0.35) {$-\pi$};
    \node[anchor=north west] at (2.6,-0.35) {$\pi$};
    \node[anchor=east] at (-1.4,1.5) {\small eğim $+$};
    \draw[->] (-1.35,1.5) -- (-0.95,1.2);
    \node[anchor=east] at (-2.3,-1.4) {\small türevi};
    \draw[->] (-2.4,-1.6) -- (-2.2,-2.1);
    % alt cizim: turev
    \begin{scope}[shift={(0,-4.1)}]
      \draw[-{Stealth[length=2.2mm]},line width=0.7pt] (-3.5,0) -- (3.8,0);
      \draw[-{Stealth[length=2.2mm]},line width=0.7pt] (0,-1.4) -- (0,1.5);
      \draw[murekkep,line width=1pt] (-2.6,0) -- (-2.6,0.9) -- (0,0.9) -- (0,-0.9) -- (2.6,-0.9) -- (2.6,0);
      \node[above] at (-1.3,0.9) {$1$};
      \node[left] at (0,-0.9) {$-1$};
      \draw[turuncukalem,line width=0.9pt] (-1.3,-0.12) -- (-1.3,0.12);
      \node[turuncukalem,below,font=\small] at (-1.3,-0.12) {$-\pi/2$};
      \draw[turuncukalem,line width=0.9pt] (1.3,-0.12) -- (1.3,0.12);
      \node[turuncukalem,below,font=\small] at (1.3,-0.12) {$\pi/2$};
    \end{scope}
  \end{tikzpicture}
\end{center}
\[
  X(f)=\frac{j}{2\pi^2}\left[\Pi\!\left(\frac{f+\pi/2}{\pi/2}\right)-\Pi\!\left(\frac{f-\pi/2}{\pi/2}\right)\right]
\]

\sayfa{3}
\[
  X(f)=\frac{j}{2\pi^2}\left[\Pi\!\left(\frac{f+\pi/2}{\pi/2}\right)-\Pi\!\left(\frac{f-\pi/2}{\pi/2}\right)\right]
\]
\[
  Y(f)=X(f)\,H(f)\qquad\left\{\ \frac{1}{10}\,\Pi\!\left(\frac{f}{40\pi}\right)\cdot e^{j2\pi f 5}=H(f)\right.
\]
\[
  Y(f)=\frac{j}{20\pi^2}\left[\Pi\!\left(\frac{f+\pi/2}{\underbrace{\pi/2}_{z}}\right)-\Pi\!\left(\frac{f-\pi/2}{\pi/2}\right)\right]\cdot e^{j2\pi f 5}
\]
\[
  Az\,\operatorname{sinc}(tz)\overset{F}{\longleftrightarrow}A\,\Pi\!\left(\frac{f}{z}\right),
  \qquad Az\rightarrow 1,\quad A\rightarrow 1,\quad z\rightarrow\frac{\pi}{2}
\]
\begin{align*}
  \frac{\pi}{2}\operatorname{sinc}\!\left(\pi\,\frac{t}{2}\right)&\donusum\Pi\!\left(\frac{f}{\pi/2}\right)\\
  \frac{\pi}{2}\operatorname{sinc}\!\left(\frac{\pi}{2}(t+5)\right)&\donusum\Pi\!\left(\frac{f}{\pi/2}\right)\cdot e^{j2\pi f 5}\\
  \frac{\pi}{2}\operatorname{sinc}\!\left(\frac{\pi}{2}(t+5)\right)\cdot e^{j2\pi\cdot\frac{\pi}{2}t}&\donusum\Pi\!\left(\frac{f+\pi/2}{\pi/2}\right)e^{j2\pi f 5}
\end{align*}
\begin{align*}
  y(t)&=\frac{j}{20\pi^2}\left[\frac{\pi}{2}\operatorname{sinc}\!\left(\frac{\pi}{2}(t+5)\right)e^{j\pi^2 t}-\frac{\pi}{2}\operatorname{sinc}\!\left(\frac{\pi}{2}(t+5)\right)e^{-j\pi^2 t}\right]\\[2mm]
  &=\frac{j}{40\pi}\operatorname{sinc}\!\left(\frac{\pi}{2}(t+5)\right)\left[e^{j\pi^2 t}-e^{-j\pi^2 t}\right]\\[2mm]
  &=-\frac{1}{20\pi}\operatorname{sinc}\!\left(\frac{\pi}{2}(t+5)\right)\cdot\underbrace{\frac{e^{j\pi^2 t}-e^{-j\pi^2 t}}{2j}}_{\sin(\pi^2 t)}\\[2mm]
  &=-\frac{1}{20\pi}\operatorname{sinc}\!\left(\frac{\pi}{2}(t+5)\right)\sin(\pi^2 t)
\end{align*}

\sayfa{4}
d)
\[
  x(t)=2\operatorname{sinc}\bigl(\underbrace{40\pi}_{z}\,t\bigr),\qquad Az=2,\quad A=\frac{2}{40\pi}
\]
\[
  \frac{1}{10}\,\Pi\!\left(\frac{f}{40\pi}\right)\cdot e^{j2\pi f 5}=H(f)
\]
\formulkutu{Az\operatorname{sinc}(tz)\overset{F}{\longleftrightarrow}A\,\Pi\!\left(\frac{f}{z}\right)}
\[ X(f)=\frac{1}{20\pi}\,\Pi\!\left(\frac{f}{40\pi}\right) \]
\[ Y(f)=X(f)\,H(f)=\frac{1}{200\pi}\,\Pi\!\left(\frac{f}{40\pi}\right)\cdot e^{j2\pi f 5} \]
\begin{align*}
  E_g(f)=|X(f)|^2&=\left(\frac{1}{200}\,\Pi\!\left(\frac{f}{40\pi}\right)\right)^{2}\\
  &=\frac{1}{40\,000}\,\Pi\!\left(\frac{f}{40\pi}\right)
\end{align*}
e) $\displaystyle E=\int_{-\infty}^{\infty}E_g(f)\,df=\int_{-20\pi}^{20\pi}\frac{1}{40000}\,df$
\[ =\frac{40\pi}{40\,000}=\frac{\pi}{1000} \]

\end{document}
